
\documentclass[11pt]{article}
\usepackage[margin=1in]{geometry}
\usepackage{amsmath,amssymb,bm}
\usepackage{siunitx}
\usepackage{hyperref}
\hypersetup{colorlinks=true,linkcolor=blue,citecolor=blue,urlcolor=blue}

\title{Apogee: Two-Stage Ascent to Circular LEO (Equatorial Launch, No Earth Rotation)}
\author{}
\date{}

\begin{document}
\maketitle

\begin{abstract}
We define and derive, from first principles, a planar (2D) rocket ascent simulator to circular low Earth orbit (LEO) aligned to publicly available Falcon~9 data. The only mission input is target circular-orbit altitude. Earth rotation is neglected by design; launch is assumed on the equator (Ecuador latitude $\approx 0^\circ$) to keep the orbit equatorial while preserving the no-rotation requirement. The model includes variable mass, non-constant gravity, U.S. Standard Atmosphere 1976 density, and aerodynamic drag. Staging, pitch-over, and cutoff are handled as events. We provide the final numbered ODE system, event logic, orbital validation, and a strict numerical methodology (adaptive Runge--Kutta + root-finding shooting).
\end{abstract}

\section{Project Definition (What We Solve)}
\subsection{Mission input}
The single mission input is the desired circular orbit altitude $h_{\text{target}}$ above mean sea level. Define
\begin{equation}
r_{\text{target}} = R_E + h_{\text{target}}.
\end{equation}

\subsection{Outputs}
The simulator must return the full time history (trajectory) until orbit insertion:
\[
t,\ r(t),\ h(t),\ \lambda(t),\ x(t),\ v(t),\ \gamma(t),\ m(t),\ q(t),\ M(t),\ D(t),
\]
plus terminal orbit diagnostics at MECO/SECO: $(a,e,r_{\text{apo}},r_{\text{peri}})$.

\subsection{Design constraints (fixed by your requirements)}
\begin{itemize}
\item \textbf{No Earth rotation} (inertial Earth is non-rotating).
\item \textbf{Equatorial launch} (Ecuador line) interpreted as: the ascent plane is the equatorial plane and the target orbit inclination is $0^\circ$. With no Earth rotation, launch latitude does not modify inertial dynamics, but we keep geometry consistent.
\item \textbf{Three guidance phases:} vertical ascent $\rightarrow$ pitch-over kick $\rightarrow$ gravity turn with thrust aligned to velocity ($\alpha=0$).
\end{itemize}

\section{Earth Constants (Required and Why)}
These constants are required because gravity and orbital mechanics are central to the problem.

\subsection{Constants}
\begin{itemize}
\item Mean Earth radius: $R_E=\SI{6371000}{m}$ (use as a geometric reference; any consistent Earth radius standard is acceptable).
\item Earth gravitational parameter: $\mu=\SI{3.986004418e14}{m^3/s^2}$ (sets $g(r)$ and orbital speed).
\item Standard gravity: $g_0=\SI{9.80665}{m/s^2}$ (definition used in $I_{sp}$ relation).
\end{itemize}

\subsection{Why these appear in the equations}
Non-constant gravity:
\begin{equation}
g(r)=\frac{\mu}{r^2}.
\end{equation}
Circular orbital speed at radius $r$:
\begin{equation}
v_{\text{circ}}(r)=\sqrt{\frac{\mu}{r}}.
\end{equation}

\section{Atmosphere Model (USSA76) and Using \texttt{ussa1976}}
\subsection{Why USSA76}
We use the \emph{U.S. Standard Atmosphere 1976} (USSA76) to obtain $\rho(h)$ and (optionally) $a_s(h)$ for Mach number. USSA76 is a published, engineering reference atmosphere used widely for trajectory and vehicle analyses \cite{NASA_USSA76}.

\subsection{Is \texttt{ussa1976} OK to use?}
Yes, as a practical implementation of USSA76: the PyPI package \texttt{ussa1976} provides a programmatic interface to the USSA76 model \cite{PyPI_USSA1976,USSA1976_Docs}. This is appropriate for APOGEE because your requirement is a consistent atmosphere model, not real-time weather.

\section{Falcon 9 Public (Official) Vehicle Inputs}
\subsection{Official Falcon 9 overview (SpaceX website)}
SpaceX lists Falcon~9 height, diameter, gross mass, and payload-to-LEO capability \cite{SpaceX_F9_Web}:
\begin{itemize}
\item Height: \SI{70}{m}
\item Core diameter: \SI{3.7}{m}
\item Vehicle mass (gross): \SI{549054}{kg}
\item Payload to LEO: \SI{22800}{kg} (capability; \textbf{payload mass is parameterized} in APOGEE)
\end{itemize}

\subsection{Official stage thrust and configuration (SpaceX User's Guide)}
SpaceX Falcon User's Guide (2025) provides stage propulsion characteristics, including engine counts and stage thrust values \cite{SpaceX_UsersGuide_2025}:
\begin{itemize}
\item Stage~1: 9 Merlin 1D (M1D), propellant LOX/RP-1; stage total thrust at sea level \SI{7686}{kN}; engine thrust per engine up to \SI{845}{kN} sea level; throttling capability noted \cite{SpaceX_UsersGuide_2025}.
\item Stage~2: 1 Merlin Vacuum (MVac), propellant LOX/RP-1; stage thrust in vacuum \SI{981}{kN}; throttling capability noted \cite{SpaceX_UsersGuide_2025}.
\end{itemize}

\subsection{Derived geometry for drag reference area}
Using diameter $d=\SI{3.7}{m}$, define reference area
\begin{equation}
A_{\text{ref}}=\pi\left(\frac{d}{2}\right)^2=\pi(1.85)^2\approx \SI{10.75}{m^2}.
\end{equation}

\subsection{What SpaceX does \emph{not} publish in these sources}
The SpaceX website and the 2025 user guide do not consistently provide full stage-by-stage \emph{dry mass}, \emph{propellant mass}, or detailed $I_{sp}$ curves. Therefore, if you want a numerically closed model \emph{today}, you must choose one of:
\begin{enumerate}
\item treat stage masses and $I_{sp}$ as \textbf{tunable engineering parameters}, calibrated to hit known staging and insertion performance; or
\item use secondary compilations (not ``official SpaceX''), clearly labeled as assumptions.
\end{enumerate}
APOGEE will proceed with option (1) by default (clean and honest).

\section{Coordinate System and State}
We use a planar, Earth-centered polar geometry (equatorial plane).

\subsection{Definitions}
\begin{itemize}
\item $r(t)$: geocentric radius
\item $h(t)=r(t)-R_E$: altitude
\item $\lambda(t)$: downrange central angle (rad)
\item $x(t)=R_E\lambda(t)$: surface arc-length downrange (m)
\item $v(t)$: speed magnitude
\item $\gamma(t)$: flight-path angle measured from local horizontal (rad)
\end{itemize}

\subsection{State vector}
\begin{equation}
\boxed{\bm{y}(t)=\begin{bmatrix} r \\ \lambda \\ v \\ \gamma \\ m \end{bmatrix}.}
\end{equation}

\section{Kinematics: Step-by-Step Derivation}
\subsection{Velocity decomposition}
In polar coordinates:
\begin{equation}
\bm{v}=\dot r\,\hat{\bm e}_r+r\dot\lambda\,\hat{\bm e}_\lambda.
\end{equation}
Define $\gamma$ such that
\begin{equation}
\dot r=v\sin\gamma,\qquad r\dot\lambda=v\cos\gamma.
\end{equation}
Hence the kinematic ODEs:
\begin{equation}
\boxed{\dot r=v\sin\gamma},\qquad
\boxed{\dot\lambda=\frac{v\cos\gamma}{r}},\qquad
\boxed{\dot x=R_E\dot\lambda=R_E\frac{v\cos\gamma}{r}}.
\end{equation}

\section{Forces and Models}
\subsection{Gravity (non-constant)}
Central gravity acceleration:
\begin{equation}
\bm{a}_g=-\frac{\mu}{r^2}\hat{\bm e}_r.
\end{equation}

\subsection{Drag (USSA76 density)}
Drag magnitude:
\begin{equation}
\boxed{D=\frac{1}{2}\rho(h)\,C_D(M)\,A_{\text{ref}}\,v^2.}
\end{equation}
Mach number:
\begin{equation}
M=\frac{v}{a_s(h)}.
\end{equation}

\subsection{Thrust direction and angle of attack}
Define unit tangent along velocity and unit normal:
\begin{equation}
\hat{\bm t}=\sin\gamma\,\hat{\bm e}_r+\cos\gamma\,\hat{\bm e}_\lambda,\qquad
\hat{\bm n}=\cos\gamma\,\hat{\bm e}_r-\sin\gamma\,\hat{\bm e}_\lambda.
\end{equation}
Let $\alpha$ be angle of attack (AoA) between thrust and velocity direction:
\begin{equation}
\hat{\bm u}_T=\cos\alpha\,\hat{\bm t}+\sin\alpha\,\hat{\bm n}.
\end{equation}
Then thrust acceleration:
\begin{equation}
\bm{a}_T=\frac{T}{m}\hat{\bm u}_T.
\end{equation}
Drag acceleration is opposite velocity:
\begin{equation}
\bm{a}_D=-\frac{D}{m}\hat{\bm t}.
\end{equation}

\section{Variable Mass: Step-by-Step Derivation}
Specific impulse definition:
\begin{equation}
I_{sp}=\frac{T}{\dot m_{\text{prop}}g_0}.
\end{equation}
Vehicle mass decreases as $\dot m=-\dot m_{\text{prop}}$, so
\[
I_{sp}=\frac{T}{(-\dot m)g_0}\quad\Rightarrow\quad
\boxed{\dot m=-\frac{T}{I_{sp}g_0}}.
\]

\section{Dynamics: Step-by-Step Derivation of $\dot v$ and $\dot\gamma$}
Total acceleration:
\begin{equation}
\bm{a}=\bm{a}_T+\bm{a}_D+\bm{a}_g.
\end{equation}

\subsection{Tangential (speed) equation}
By definition of speed along the velocity direction:
\begin{equation}
\dot v=\bm{a}\cdot\hat{\bm t}.
\end{equation}
Compute each term:
\[
\bm{a}_T\cdot\hat{\bm t}=\frac{T}{m}\hat{\bm u}_T\cdot\hat{\bm t}=\frac{T}{m}\cos\alpha,
\quad
\bm{a}_D\cdot\hat{\bm t}=-\frac{D}{m},
\quad
\bm{a}_g\cdot\hat{\bm t}=\left(-\frac{\mu}{r^2}\hat{\bm e}_r\right)\cdot\hat{\bm t}=-\frac{\mu}{r^2}\sin\gamma.
\]
Thus
\begin{equation}
\boxed{\dot v=\frac{T}{m}\cos\alpha-\frac{D}{m}-\frac{\mu}{r^2}\sin\gamma.}
\end{equation}

\subsection{Normal (turning) equation}
A standard polar-kinematics identity yields
\begin{equation}
\bm{a}\cdot\hat{\bm n}=v\dot\gamma-\frac{v^2}{r}\cos\gamma.
\end{equation}
Now compute $\bm{a}\cdot\hat{\bm n}$ from forces:
\[
\bm{a}_T\cdot\hat{\bm n}=\frac{T}{m}\sin\alpha,\qquad
\bm{a}_D\cdot\hat{\bm n}=0,\qquad
\bm{a}_g\cdot\hat{\bm n}=\left(-\frac{\mu}{r^2}\hat{\bm e}_r\right)\cdot\hat{\bm n}=-\frac{\mu}{r^2}\cos\gamma.
\]
So
\[
\frac{T}{m}\sin\alpha-\frac{\mu}{r^2}\cos\gamma=v\dot\gamma-\frac{v^2}{r}\cos\gamma.
\]
Solve for $\dot\gamma$:
\begin{equation}
\boxed{\dot\gamma=\frac{T}{mv}\sin\alpha+\left(\frac{v}{r}-\frac{\mu}{r^2v}\right)\cos\gamma.}
\end{equation}

\section{Gravity Turn Justification and Simplification}
In the gravity-turn phase we enforce \textbf{thrust aligned with velocity}:
\begin{equation}
\boxed{\alpha=0\quad\Rightarrow\quad \bm{T}\parallel\bm{V}.}
\end{equation}
Then the coupled equations simplify to
\begin{equation}
\boxed{\dot v=\frac{T}{m}-\frac{D}{m}-\frac{\mu}{r^2}\sin\gamma},
\qquad
\boxed{\dot\gamma=\left(\frac{v}{r}-\frac{\mu}{r^2v}\right)\cos\gamma}.
\end{equation}
This is physically consistent: thrust produces tangential acceleration, while gravity and curvature produce turning.

\section{Final Numbered ODE System (Implementation-Ready)}
Define closures:
\[
h=r-R_E,\quad \rho=\rho(h)\ \text{(USSA76)},\quad a_s=a_s(h)\ \text{(USSA76)},\quad M=v/a_s,
\]
\[
D=\frac{1}{2}\rho(h)\,C_D(M)\,A_{\text{ref}}\,v^2.
\]
Then
\begin{align}
\boxed{\dot r} &= \boxed{v\sin\gamma} \label{eq:r}\\
\boxed{\dot\lambda} &= \boxed{\frac{v\cos\gamma}{r}} \label{eq:lambda}\\
\boxed{\dot v} &= \boxed{\frac{T}{m}\cos\alpha-\frac{D}{m}-\frac{\mu}{r^2}\sin\gamma} \label{eq:v}\\
\boxed{\dot\gamma} &= \boxed{\frac{T}{mv}\sin\alpha+\left(\frac{v}{r}-\frac{\mu}{r^2v}\right)\cos\gamma} \label{eq:gamma}\\
\boxed{\dot m} &= \boxed{-\frac{T}{I_{sp}g_0}} \label{eq:m}
\end{align}
During coast: set $T=0$ so $\dot m=0$.

\section{Hybrid Events (Falcon 9--Aligned, Simplified)}
\subsection{Phase A: Vertical ascent}
Hold $\gamma=\pi/2$ until a trigger (choose altitude $h_{po}$).
Event:
\[
g_{po}(\bm{y}) = (r-R_E) - h_{po} = 0.
\]
Purpose: avoid the $1/v$ term near $v\to 0$ in \eqref{eq:gamma}.

\subsection{Phase B: Pitch-over kick}
At $t=t_{po}$ set (instantaneous reset):
\begin{equation}
\boxed{\gamma(t_{po}^+)=\frac{\pi}{2}-\theta_0.}
\end{equation}

\subsection{Phase C: Gravity turn}
Set $\alpha=0$ for the remainder of powered ascent.

\subsection{Stage 1 $\rightarrow$ Stage 2 separation (event-driven)}
We model a two-stage stack consistent with Falcon~9 being a two-stage LOX/RP-1 vehicle \cite{SpaceX_UsersGuide_2025}.
In the simplified model:
\begin{itemize}
\item Stage~1 burn ends when $m$ reaches a stage-1 burnout mass $m_{1,\text{end}}$.
\item Then drop stage-1 dry mass instantly and switch $(T,I_{sp},C_D,A)$ to stage-2 values.
\end{itemize}
Event (burnout):
\[
g_{\text{S1}}(\bm{y})=m-m_{1,\text{end}}=0.
\]
Reset:
\[
m^+ = m^- - m_{1,\text{dry}}.
\]

\subsection{Final cutoff (MECO/SECO) chosen for convergence}
To guarantee a solvable shooting problem, use a cutoff mass $m_{\text{cut}}$ on the final stage:
\[
g_{\text{cut}}(\bm{y}) = m - m_{\text{cut}} = 0,\qquad T\gets 0.
\]

\section{Orbit Validation (Two-Body Mechanics)}
At cutoff, compute
\begin{equation}
\varepsilon=\frac{v^2}{2}-\frac{\mu}{r},\qquad
h_{\text{ang}}=rv\cos\gamma.
\end{equation}
Eccentricity:
\begin{equation}
e=\sqrt{1+\frac{2\varepsilon h_{\text{ang}}^2}{\mu^2}}.
\end{equation}
Semi-major axis (bound orbit: $\varepsilon<0$):
\begin{equation}
a=-\frac{\mu}{2\varepsilon}.
\end{equation}
Apoapsis/periapsis radii:
\begin{equation}
r_{\text{apo}}=a(1+e),\qquad r_{\text{peri}}=a(1-e).
\end{equation}

\section{Shooting Method (What We Solve Numerically)}
\subsection{Decision variables}
We solve for
\[
\bm{u}=\begin{bmatrix}\theta_0\\ m_{\text{cut}}\end{bmatrix}.
\]

\subsection{Residual equations}
We target circular insertion at $r_{\text{target}}$:
\begin{equation}
F_1(\bm{u})=r_{\text{cut}}(\bm{u})-r_{\text{target}},
\qquad
F_2(\bm{u})=v_{\text{cut}}(\bm{u})-\sqrt{\frac{\mu}{r_{\text{target}}}}.
\end{equation}
Success diagnostics (not necessarily enforced as equalities at first):
\[
|\gamma_{\text{cut}}|\le \epsilon_\gamma,\qquad |r_{\text{apo}}-r_{\text{target}}|\le \epsilon_r,\qquad |r_{\text{peri}}-r_{\text{target}}|\le \epsilon_r.
\]

\section{Numerical Methods (Strict Recommendations)}
\subsection{ODE integration}
Use adaptive embedded explicit Runge--Kutta (nonstiff):
\begin{itemize}
\item Dormand--Prince 5(4) (Dopri5), or Tsitouras 5(4) (Tsit5).
\end{itemize}
Use event detection (root finding on $g(\bm{y})$) for: pitch-over, stage separation, cutoff.

\subsection{Nonlinear solve for shooting}
Solve $\bm{F}(\bm{u})=\bm{0}$ using:
\begin{itemize}
\item Powell hybrid / trust-region (robust) with finite-difference Jacobian, or
\item Broyden (quasi-Newton) for fewer trajectory evaluations.
\end{itemize}
Scale residuals:
\[
\tilde F_1=\frac{F_1}{r_{\text{target}}},\qquad \tilde F_2=\frac{F_2}{v_{\text{circ}}(r_{\text{target}})}.
\]

\section{Final ``Start Coding'' Input Set (Official + Explicit Assumptions)}
\subsection{Official inputs (from SpaceX)}
\begin{itemize}
\item Falcon~9 diameter $d=\SI{3.7}{m}$ and gross mass $\SI{549054}{kg}$ \cite{SpaceX_F9_Web}.
\item Stage thrust: Stage~1 sea-level stage thrust $\SI{7686}{kN}$ and Stage~2 vacuum thrust $\SI{981}{kN}$ \cite{SpaceX_UsersGuide_2025}.
\item Propellant: LOX/RP-1 for both stages \cite{SpaceX_UsersGuide_2025}.
\end{itemize}

\subsection{Atmosphere (official standard)}
\begin{itemize}
\item USSA76: $\rho(h)$ and $a_s(h)$ from NASA USSA76 \cite{NASA_USSA76}, implemented via \texttt{ussa1976} \cite{PyPI_USSA1976,USSA1976_Docs}.
\end{itemize}

\subsection{Assumptions you must choose (because SpaceX does not publish them here)}
These are \emph{required} to numerically integrate \eqref{eq:m} and to time staging by mass depletion:
\begin{itemize}
\item $I_{sp,1}$, $I_{sp,2}$ (use constant per stage in v1).
\item Stage dry masses and propellant masses (or, equivalently, $m_{1,\text{end}}$ and $m_{\text{cut}}$ bounds).
\item Drag model $C_D(M)$ (we recommend a smooth piecewise function; tune later).
\end{itemize}

\textbf{Minimal viable} (to run immediately):
\begin{itemize}
\item Choose $I_{sp,1}$ and $I_{sp,2}$ as constants; choose stage mass split consistent with total gross mass.
\item Choose $C_D(M)$ simple and stable (piecewise + smoothing), and accept it as an uncertainty knob.
\end{itemize}
%===========================================================
% APOGEE — Calibration Section (Falcon 9-aligned, strict math)
% Paste this into your master LaTeX.
%===========================================================

\section{Falcon 9--Aligned Calibration (Closing the Model Without Guessing)}

\subsection{What is unknown and why calibration is necessary}
SpaceX public sources provide (reliably) vehicle geometry and thrust levels, plus some timeline-type quantities (e.g., second-stage burn time on the Falcon 9 web page). However, these public sources do not fully specify a complete, internally consistent set of \emph{stage propellant masses}, \emph{stage dry masses}, and \emph{$I_{sp}$ curves} required to integrate the variable-mass ODE system with staging \cite{SpaceX_F9_Web,SpaceX_UsersGuide_2025}. 

Therefore, APOGEE closes the system using a \textbf{calibration problem}:
we choose the \emph{smallest set} of missing constants and solve for them using (i) exact physical identities and (ii) known Falcon 9 flight-consistent constraints. This is the standard engineering approach for early-phase ascent simulation when proprietary vehicle databases are unavailable.

\subsection{Known (official) inputs used as hard constraints}
We fix the following as non-negotiable inputs from SpaceX:
\begin{itemize}
\item Core diameter $d=\SI{3.7}{m}$ and gross liftoff mass $m_0=\SI{549054}{kg}$ \cite{SpaceX_F9_Web}.
\item Stage thrust (stage totals): Stage~1 $T_{1,\text{SL}}=\SI{7686}{kN}$ (sea level) and Stage~2 $T_{2,\text{vac}}=\SI{981}{kN}$ (vacuum) \cite{SpaceX_UsersGuide_2025}.
\item Stage~2 burn time shown on the Falcon 9 vehicle page: $t_{2,\text{burn}}=\SI{397}{s}$ \cite{SpaceX_F9_Web}.
\item Propellant type: LOX/RP-1 for both stages \cite{SpaceX_UsersGuide_2025}.
\item Atmosphere: U.S. Standard Atmosphere 1976 (USSA76), implemented via \texttt{ussa1976}, which explicitly implements NASA TM-X-74335 \cite{NASA_USSA76,PyPI_USSA1976}.
\end{itemize}

\subsection{State and equations used during calibration}
We calibrate using the same implementation-ready ODEs (Eqs.~\eqref{eq:r}--\eqref{eq:m} in the main document):
\[
\dot r=v\sin\gamma,\quad
\dot\lambda=\frac{v\cos\gamma}{r},\quad
\dot v=\frac{T}{m}\cos\alpha-\frac{D}{m}-\frac{\mu}{r^2}\sin\gamma,\quad
\dot\gamma=\frac{T}{mv}\sin\alpha+\left(\frac{v}{r}-\frac{\mu}{r^2v}\right)\cos\gamma,\quad
\dot m=-\frac{T}{I_{sp}g_0}.
\]
During gravity-turn guidance we enforce $\alpha=0$ (thrust aligned with velocity). Drag is
\[
D=\frac{1}{2}\rho(h)\,C_D(M)\,A_{\text{ref}}\,v^2,\quad
A_{\text{ref}}=\pi\left(\frac{d}{2}\right)^2,
\quad M=\frac{v}{a_s(h)}.
\]

\subsection{Calibration variables (minimal set)}
To obtain a numerically closed model that is Falcon 9-aligned but still simple, we define the \emph{calibration parameter vector}
\begin{equation}
\bm{p}=
\begin{bmatrix}
I_{sp,1} \\
I_{sp,2} \\
t_{1,\text{burn}} \\
m_{1,\text{dry}} \\
m_{2,\text{dry}} \\
\bm{c}_{D}
\end{bmatrix},
\end{equation}
where $\bm{c}_D$ are the few parameters defining a smooth $C_D(M)$ model (see Sec.~\ref{sec:cdmodel}).

\paragraph{Why include $t_{1,\text{burn}}$?}
The Falcon 9 web page explicitly lists Stage~2 burn time \cite{SpaceX_F9_Web}; Stage~1 burn time is not consistently presented with the same clarity in those same official sources. To remain strict and avoid mixing inconsistent datasets, we treat $t_{1,\text{burn}}$ as a calibration variable constrained to the physically plausible range of Stage~1 burn duration (close to the commonly cited $\sim\SI{160}{s}$). If you later confirm a Stage~1 burn time from an official source in your docs, lock it as a constant.

\subsection{Stage propellant masses implied by thrust, burn time, and $I_{sp}$}
The variable-mass identity
\[
\dot m=-\frac{T}{I_{sp}g_0}
\]
implies that for a stage with constant $(T,I_{sp})$ during its burn, the propellant mass consumed is
\begin{equation}\label{eq:mprop}
m_{i,\text{prop}}( \bm{p} )=\int_0^{t_{i,\text{burn}}}\!\!\frac{T_i}{I_{sp,i}g_0}\,dt
=\frac{T_i\,t_{i,\text{burn}}}{I_{sp,i}\,g_0}.
\end{equation}
This is \textbf{not an assumption}; it follows directly from the definition of $I_{sp}$ and mass conservation.

\subsection{Mass accounting constraint (hard equality)}
At liftoff, the gross mass must equal the sum of dry masses, propellant masses, and payload:
\begin{equation}\label{eq:massbudget}
m_0=
m_{1,\text{dry}} + m_{2,\text{dry}}
+ m_{1,\text{prop}}(\bm{p}) + m_{2,\text{prop}}(\bm{p})
+ m_{\text{PL}},
\end{equation}
where $m_{\text{PL}}$ is \textbf{parameterized} (mission input). Equation \eqref{eq:massbudget} is enforced as a strict constraint in calibration.

\subsection{Falcon 9 flight-consistency constraints (targets)}
We do \emph{not} require reproducing a specific mission; instead we enforce a small set of flight-consistency constraints that any Falcon 9-like ascent must satisfy:

\paragraph{(1) Orbit insertion conditions (hard targets).}
At final cutoff (SECO), for a circular orbit at $r_{\text{target}}=R_E+h_{\text{target}}$:
\begin{equation}\label{eq:orbit_targets}
r_{\text{SECO}} = r_{\text{target}},\qquad
v_{\text{SECO}} = \sqrt{\frac{\mu}{r_{\text{target}}}},\qquad
\gamma_{\text{SECO}}\approx 0.
\end{equation}
(We typically enforce the first two as hard equalities via shooting; the third as a small tolerance.)

\paragraph{(2) Stage event ordering (hard logic).}
Stage~1 burn $\rightarrow$ MECO $\rightarrow$ separation $\rightarrow$ Stage~2 ignition $\rightarrow$ SECO.
This is implemented via event functions on $m(t)$ and time.

\paragraph{(3) Dynamic pressure sanity (soft constraint).}
Maximum dynamic pressure
\[
q(t)=\frac{1}{2}\rho(h(t))v(t)^2
\]
must occur in the lower atmosphere and remain within a plausible range. Since SpaceX throttle schedules are not modeled, we treat this as a diagnostic or a soft penalty to prevent unrealistic drag settings. (We do \emph{not} force a specific Max-Q value unless you decide to.)

\subsection{Objective function (strict, dimensionless, well-conditioned)}
Define residuals for a single simulation run:
\begin{align}
F_1(\bm{p},\bm{u})&=\frac{r_{\text{SECO}}(\bm{p},\bm{u})-r_{\text{target}}}{r_{\text{target}}},\\
F_2(\bm{p},\bm{u})&=\frac{v_{\text{SECO}}(\bm{p},\bm{u})-\sqrt{\mu/r_{\text{target}}}}{\sqrt{\mu/r_{\text{target}}}},\\
F_3(\bm{p},\bm{u})&=\frac{\gamma_{\text{SECO}}(\bm{p},\bm{u})}{\gamma_{\text{scale}}},\\
F_4(\bm{p})&=\frac{m_0-\left(m_{1,\text{dry}}+m_{2,\text{dry}}+m_{1,\text{prop}}(\bm{p})+m_{2,\text{prop}}(\bm{p})+m_{\text{PL}}\right)}{m_0},
\end{align}
where $\gamma_{\text{scale}}$ is a small scaling (e.g., $\gamma_{\text{scale}}=1^\circ$ in radians).

We split unknowns into:
\begin{itemize}
\item \textbf{Guidance/shooting variables} $\bm{u}$ (solved per mission): 
$\bm{u}=[\theta_0,\ m_{\text{cut}}]^T$.
\item \textbf{Vehicle calibration parameters} $\bm{p}$ (solved once, reused).
\end{itemize}

The calibration solves:
\begin{equation}\label{eq:calib_problem}
\min_{\bm{p}}\; 
\left\|\bm{F}(\bm{p},\bm{u}^\star(\bm{p}))\right\|_2^2
+\beta\,\Phi(\bm{p}),
\end{equation}
subject to physical bounds (positivity, plausible $I_{sp}$ ranges, etc.), where:
\begin{itemize}
\item $\bm{u}^\star(\bm{p})$ is obtained by an \textbf{inner shooting solve} enforcing \eqref{eq:orbit_targets} for a reference mission (e.g., $h_{\text{target}}=\SI{200}{km}$, chosen once).
\item $\Phi(\bm{p})$ is a \textbf{regularizer} encoding ``do not choose absurd values'', e.g.,
\[
\Phi(\bm{p})=
\left(\frac{I_{sp,1}-\bar I_{sp,1}}{\sigma_{I1}}\right)^2+
\left(\frac{I_{sp,2}-\bar I_{sp,2}}{\sigma_{I2}}\right)^2+
\left(\frac{t_{1,\text{burn}}-\bar t_1}{\sigma_{t1}}\right)^2,
\]
with $\bar I_{sp,i}$ and $\bar t_1$ chosen as nominal engineering values.
\end{itemize}

\paragraph{Strictness note.}
Equation \eqref{eq:mprop} and mass budget \eqref{eq:massbudget} ensure full dimensional consistency and remove degrees of freedom. The regularizer $\Phi$ is not ``cheating'': it is the formal way to encode prior engineering knowledge when public data is incomplete.

\subsection{Recommended numerical solution strategy (robust)}
\subsubsection{Inner loop: mission shooting}
Given $\bm{p}$, solve for $\bm{u}=[\theta_0,m_{\text{cut}}]$ by enforcing two equalities:
\[
r_{\text{SECO}}(\bm{p},\bm{u})-r_{\text{target}}=0,\qquad
v_{\text{SECO}}(\bm{p},\bm{u})-\sqrt{\mu/r_{\text{target}}}=0.
\]
Use a robust root-finder (Powell hybrid / trust-region) with finite-difference Jacobian. Each residual evaluation integrates the ODEs with event handling.

\subsubsection{Outer loop: calibration}
With $\bm{u}^\star(\bm{p})$ from the inner loop, minimize \eqref{eq:calib_problem} using a derivative-free trust-region method (or quasi-Newton if you later move to autodiff). Because each function evaluation is expensive (full trajectory), keep $\bm{p}$ low-dimensional.

\subsubsection{ODE integration}
Use adaptive embedded Runge--Kutta 5(4) (Tsit5 or Dopri5) with event detection. This is appropriate because the system is nonlinear but not stiff in typical ascent regimes.

\subsection{Aerodynamic model to calibrate (simple, stable)}
\label{sec:cdmodel}
SpaceX does not publish Falcon 9 $C_D(M)$; thus we use a smooth function with few parameters:
\begin{equation}\label{eq:cd_smooth}
C_D(M)=C_{D,\text{sub}} 
+ \Delta C_{D,\text{tr}}\;S(M;M_1,M_2)
+ \Delta C_{D,\text{sup}}\;S(M;M_3,M_4),
\end{equation}
where $S$ is a smooth step (e.g., logistic or $\tanh$-based), and $(M_1<M_2<M_3<M_4)$ define transition ranges (subsonic$\rightarrow$transonic$\rightarrow$supersonic). The parameter vector $\bm{c}_D$ consists of these amplitudes and transition Mach numbers. This ensures differentiability and numerical stability.

\subsection{Initial conditions and ``Ecuador launch'' (with no Earth rotation)}
With Earth rotation explicitly neglected, the only consistent interpretation of ``Ecuador launch'' is that the trajectory is modeled in the equatorial plane and the target orbit is equatorial. The initial condition is:
\[
r(0)=R_E,\quad \lambda(0)=0,\quad v(0)=0,\quad \gamma(0)=\frac{\pi}{2},\quad m(0)=m_0.
\]
No extra inertial boost is applied.

\subsection{Calibration deliverable (what you hand to coding AIs)}
The calibration outputs a fully specified parameter set:
\[
\boxed{
\left\{
I_{sp,1},\ I_{sp,2},\ t_{1,\text{burn}},\ 
m_{1,\text{dry}},\ m_{2,\text{dry}},\ 
C_D(M)\ \text{parameters}
\right\}
}
\]
together with the fixed official constants $(d,m_0,T_{1,\text{SL}},T_{2,\text{vac}},t_{2,\text{burn}})$ and atmosphere model choice (USSA76 via \texttt{ussa1976}).

\subsection{Key strict validation checks (must pass)}
Any calibrated solution is rejected unless:
\begin{itemize}
\item Mass budget \eqref{eq:massbudget} holds to numerical tolerance ($<10^{-8}$ relative).
\item $m_{i,\text{prop}}>0$, $m_{i,\text{dry}}>0$.
\item Orbit diagnostics satisfy near-circularity at insertion:
$|e|<e_{\max}$ (e.g. $10^{-3}$) and $|r_{\text{apo}}-r_{\text{peri}}|$ small compared to $r_{\text{target}}$.
\item Dynamic pressure $q(t)$ remains physically plausible and peaks in the lower atmosphere (diagnostic for $C_D$ sanity).
\end{itemize}

%===========================================================
% IEEE citations used above (already in your main bibliography)
%===========================================================
% \cite{SpaceX_F9_Web,SpaceX_UsersGuide_2025,NASA_USSA76,PyPI_USSA1976}

\section{IEEE References}
\bibliographystyle{IEEEtran}
\begin{thebibliography}{99}

\bibitem{SpaceX_UsersGuide_2025}
SpaceX, ``Falcon User's Guide,'' May 9, 2025. [Online]. Available: \url{https://www.spacex.com/assets/media/falcon-users-guide-2025-05-09.pdf}. Accessed: 2026-01-10.

\bibitem{SpaceX_F9_Web}
SpaceX, ``Falcon 9,'' (vehicle overview page). [Online]. Available: \url{https://www.spacex.com/vehicles/falcon-9/}. Accessed: 2026-01-10.

\bibitem{NASA_USSA76}
NASA, ``U.S. Standard Atmosphere, 1976 (NASA-TM-X-74335),'' 1976. [Online]. Available: \url{https://ntrs.nasa.gov/api/citations/19770009539/downloads/19770009539.pdf}. Accessed: 2026-01-10.

\bibitem{PyPI_USSA1976}
PyPI, ``ussa1976.'' [Online]. Available: \url{https://pypi.org/project/ussa1976/}. Accessed: 2026-01-10.

\bibitem{USSA1976_Docs}
USSA1976 Documentation, ``The U.S. Standard Atmosphere 1976 model.'' [Online]. Available: \url{https://ussa1976.readthedocs.io/}. Accessed: 2026-01-10.

\end{thebibliography}

\end{document}
